\providecommand{\kBKM}{\kappa_{\mathrm{BKM}}}
\providecommand{\kch}{\kappa_{\mathrm{ch}}}
\providecommand{\kcat}{\kappa_{\mathrm{cat}}}
\providecommand{\kfib}{\kappa_{\mathrm{fiber}}}
\providecommand{\kHeis}{\kappa_{\mathrm{ch}}^{\mathrm{Heis}}}
\providecommand{\Z}{\mathbb{Z}}
\providecommand{\C}{\mathbb{C}}

\section*{\texorpdfstring{$N=4$}{N=4} Primitive CHL Denominator Data}
\label{wn:sec:N4-primitive-chl-data}

\subsection*{Primitive CHL Normalisation}

The primitive CHL denominator branch attached to the cyclic orders
\[
  N \in \{1,2,3,4,6\}
\]
uses the \(g_N\)-twined K3 Jacobi input in the
Gritsenko-Nikulin additive lift.  Borcherds' singular theta weight
formula gives
\[
  \kBKM(\Phi_N)
  =
  \operatorname{wt}(\Phi_N)
  =
  \frac{c_N(0)}{2}.
\]
The constants are
\[
\begin{array}{c|c|c|c}
N & \text{Mathieu frame shape} & c_N(0) & \kBKM(\Phi_N)\\
\hline
1 & 1^{24} & 10 & 5\\
2 & 1^8 2^8 & 8 & 4\\
3 & 1^6 3^6 & 6 & 3\\
4 & 1^4 2^2 4^4 & 4 & 2\\
6 & 1^2 2^2 3^2 6^2 & 2 & 1
\end{array}
\]
Thus the order-four primitive denominator satisfies
\[
  c_4(0)=4,\qquad
  \kBKM(\Phi_4)=\frac{c_4(0)}{2}=2.
\]

\subsection*{Order-Four Mathieu Class}

The order-four primitive CHL branch uses the \(M_{24}\) class \(4B\).
The ordinary \(M_{24}\) twining table has \(26\) conjugacy classes.  Its
order-four frame-shape entries are
\[
\begin{array}{c|c|c}
[g] & \text{frame shape on }24\text{ letters} & c^{(g)}(0,0)\\
\hline
4A & 2^4 4^4 & 0\\
4B & 1^4 2^2 4^4 & 4\\
4C & 4^6 & -4
\end{array}
\]
and the full ordinary class-constant table is
\[
\begin{array}{c|rrrrrrrrrrrrr}
[g] &1A&2A&2B&3A&3B&4A&4B&4C&5A&6A&6B&7A&7B\\
\hline
c^{(g)}(0,0)&20&4&-4&2&-4&0&4&-4&0&0&0&-1&-1
\end{array}
\]
\[
\begin{array}{c|rrrrrrrrrrrrr}
[g] &8A&10A&11A&12A&12B&14A&14B&15A&15B&21A&21B&23A&23B\\
\hline
c^{(g)}(0,0)&0&-4&-2&0&0&0&0&0&0&-1&-1&0&0
\end{array}
\]
The \(4B\) frame shape has \(4+2\cdot 2+4\cdot 4=24\) letters and
discriminant-zero constant \(4\), giving the primitive weight \(2\).

\subsection*{Jacobi Basis}

The weak Jacobi forms of weight \(0\) and index \(1\) are generated over
modular forms by the Eichler-Zagier pair
\[
  A=\frac{1}{4}\phi_{0,1}(\tau,z),
  \qquad
  B=\phi_{-2,1}(\tau,z)
   =\frac{\nu_{1,1}(\tau,z)^2}{\eta(\tau)^6}.
\]
For a commuting cyclic pair \((g_4,h_4^s)\), \(s\in\Z/4\Z\), the
twisted-twined genus lies in the corresponding weak-Jacobi space and
has an \(A,B\)-expansion
\[
  \phi^{(g_4,h_4^s)}_{0,1}(\tau,z)
  =
  \alpha_4(s)\phi_{0,1}(\tau,z)
  +
  \beta_4(s)\phi_{-2,1}(\tau,z).
\]
The residue \(s=0\) supplies the primitive Borcherds constant
\(c_4(0)=4\).  The residues \(s=1,2,3\) supply multiplier and
centraliser data for twisted-twined refinements.

\subsection*{Gritsenko-Clery Diagonal-Divisor Catalogue}

The Gritsenko-Clery diagonal-divisor catalogue is a separate
eight-row dd-modular classification.  Its triples \((t,N;k)\) and
twice-weight constants are
\[
\begin{array}{c|c|c}
(t,N;k) & 2k & \text{form}\\
\hline
(1,1;5) & 10 & \Delta_5\\
(2,1;2) & 4 & \Delta_2\\
(1,2;3) & 6 & \nabla_3\\
(3,1;1) & 2 & \Delta_1\\
(1,3;2) & 4 & \nabla_2\\
(4,1;\tfrac12) & 1 & \Delta_{1/2}\\
(1,4;\tfrac32) & 3 & \nabla_{3/2}\\
(2,2;1) & 2 & Q_1
\end{array}
\]
Its twice-weight tuple is
\[
  (10,4,6,2,4,1,3,2),
\]
while the primitive CHL constant tuple is
\[
  (10,8,6,4,2).
\]
The intersection of the two lists is the row \((t,N;k)=(1,1;5)\), where
\(\Delta_5\) has weight \(5\).  The \(N=4\) primitive CHL entry has
\((c_4(0),\kBKM)=(4,2)\); the Gritsenko-Clery row \((1,4;\tfrac32)\)
has twice-weight constant \(3\) and weight \(3/2\).

\subsection*{\texorpdfstring{$K3\times E$}{K3 x E} Comparison Row}

For \(Y=K3\times E\), the compact Hodge column is
\[
  (h^{0,0},h^{0,1},h^{0,2},h^{0,3})=(1,1,1,1),
\]
so the compact chiral Hodge supertrace is
\[
  \kch(Y)=\sum_q(-1)^qh^{0,q}(Y)=1-1+1-1=0.
\]
K\"unneth gives
\[
  \kcat(Y)=\chi(\mathcal O_{K3})\chi(\mathcal O_E)=2\cdot0=0.
\]
The comparison row on \(K3\times E\) is
\[
\begin{array}{c|c|l}
\text{construction} & \text{value} & \text{source of the scalar}\\
\hline
\kcat(Y) & 0 & \chi(\mathcal O_{K3})\chi(\mathcal O_E)\\
\kch(Y) & 0 & \text{compact Hodge supertrace}\\
\kHeis(Y) & 3 & \text{Heisenberg-Mukai specialisation}\\
\kBKM(\Delta_5) & 5 & c_1(0)/2=10/2=\operatorname{wt}(\Delta_5)\\
\kfib(Y) & 24 & \operatorname{rk}\widetilde H(K3,\Z)
\end{array}
\]
Thus the visible spectrum of named branches is
\[
  \{\kcat,\kHeis,\kBKM,\kfib\}(Y)=\{0,3,5,24\},
\]
with the compact value \(\kch(Y)=0\) recorded separately from the
Heisenberg specialisation.

\subsection*{Rank-Three Boundary}

The formula \(\kBKM(\Phi_N)=c_N(0)/2\) fixes the automorphic weight of
the primitive denominator.  A rank-three BKM presentation at \(N=4\)
requires the integral Lorentzian root lattice, primitive real roots,
Weyl vector, imaginary-root cone, parity rule, multiplicity function,
and denominator identity.  The preceding tables record the proved
weight and class constants.

\paragraph{Sources.}
Borcherds 1998, Theorem 13.3; Gritsenko and Nikulin 1995, Part II,
Theorem 2.1; Gritsenko 1999, Corollary 3.3; Gritsenko and Clery 2008,
Theorems 1.4 and 3.1; Clery, Gritsenko and Hulek 2015 for the \(N=6\)
paramodular row; Eichler and Zagier 1985, Theorem 9.3; Eguchi, Ooguri
and Tachikawa 2011; Cheng, Harrison, Paquette and Volpato 2014;
Cheng, Duncan and Harvey 2014; Mukai 1988; Nikulin 1979.
